% !TEX root = ../main.tex

\chapter{相关工作}

鞋履设计草图到三维模型的重建任务位于二维图像理解、三维形状生成、条件扩散生成和领域适配等多个研究方向的交叉位置。本章围绕本文系统路线，分别回顾单图像三维重建、基于扩散先验的三维生成、多视图生成与前馈式三维重建、草图条件控制与领域适配等相关工作，并分析这些方法对本文任务的启发和局限。图\ref{fig:related-work-map}概括了相关技术路线与本文任务之间的关系。

\begin{figure}[htbp]
  \centering
  \includegraphics[width=0.95\textwidth]{related_work/related_work_map.png}
  \caption{相关技术路线与本文鞋履草图三维重建任务的关系}
  \label{fig:related-work-map}
\end{figure}

\section{单图像三维重建}

单图像三维重建旨在从一张二维图像恢复目标物体的三维几何结构。该问题本质上是不适定的，因为单张图像只能观测到物体的一部分表面，深度、背面和内部结构都需要模型根据先验进行推断。若将输入图像记为$I$，三维形状记为$S$，该问题可以抽象为条件概率估计：
\begin{equation}
  \hat{S} = \arg\max_S p(S \mid I).
  \label{eq:image-to-3d-map}
\end{equation}
由于从$I$到$S$存在多解性，实际方法通常需要引入类别形状先验、几何正则项或多视图约束。尽管如此，单图像重建由于输入成本低、交互简单，在产品预览、三维资产生成和快速原型设计中具有重要应用价值。

早期深度学习方法通常将二维图像特征映射到体素、点云或网格等三维表示。Pixel2Mesh使用图卷积网络从椭球网格逐步变形得到目标物体网格，体现了从图像特征直接预测三维网格顶点的思路\citep{wang2018pixel2mesh}。Mesh R-CNN进一步将实例分割、体素预测和网格细化结合起来，在检测目标的同时生成三维形状\citep{gkioxari2019meshrcnn}。这类方法为图像到网格重建提供了较早的端到端范式，但它们通常依赖特定类别数据集训练，面对鞋履设计草图这类风格化输入时泛化能力有限。

另一类重要工作是神经隐式表示和神经辐射场。NeRF通过多层感知机表示空间位置和观察方向到颜色、密度的映射，在多视角图像条件下实现了高质量新视角合成\citep{mildenhall2020nerf}。其核心体渲染思想可写为
\begin{equation}
  C(\mathbf{r}) =
  \int_{t_n}^{t_f} T(t)\sigma(\mathbf{r}(t))\mathbf{c}(\mathbf{r}(t), \mathbf{d})\,dt,
  \quad
  T(t)=\exp\left(-\int_{t_n}^{t}\sigma(\mathbf{r}(s))\,ds\right),
  \label{eq:nerf-rendering}
\end{equation}
其中$\mathbf{r}(t)$表示相机射线，$\sigma$表示体密度，$\mathbf{c}$表示颜色，$T(t)$表示从近端到$t$位置的透射率。NeRF的优势在于连续空间表达和视角一致性，但经典NeRF通常需要同一场景的多视角图像和较长优化时间。对于本文关注的单视图或三视图设计草图输入而言，NeRF类方法能够提供三维表示思想，但直接用于本地快速鞋履重建仍存在输入条件和效率限制。

总体来看，单图像三维重建方法证明了从二维视觉信息预测三维形状的可行性，但也暴露出本文任务中的关键难点：鞋履设计草图缺少真实纹理和光照，且鞋履内部空间、鞋底层次和鞋面厚度等不可见结构无法由单视图直接确定。因此，本文没有采用从零训练的类别专用重建网络，而是利用开源图像到三维生成模型作为基础能力，并通过草图预处理和条件生成模块减小输入域差异。

\section{基于扩散先验的三维生成}

扩散模型在二维图像生成领域表现出强大的先验建模能力。典型扩散模型首先在前向过程中逐步向真实图像$x_0$加入高斯噪声：
\begin{equation}
  x_t = \sqrt{\bar{\alpha}_t}x_0 + \sqrt{1-\bar{\alpha}_t}\epsilon,
  \quad \epsilon \sim \mathcal{N}(0, I),
  \label{eq:diffusion-forward}
\end{equation}
然后训练网络$\epsilon_\theta$预测噪声，其常用目标函数为
\begin{equation}
  \mathcal{L}_{\mathrm{diff}} =
  \mathbb{E}_{x_0,t,\epsilon}
  \left[
    \left\|\epsilon-\epsilon_\theta(x_t,t,c)\right\|_2^2
  \right],
  \label{eq:diffusion-loss}
\end{equation}
其中$c$可以表示文本、图像或控制图等条件。随着文本到图像和图像到图像扩散模型的发展，研究者开始探索如何将二维扩散先验迁移到三维生成任务中。DreamFusion是这一方向的代表性工作之一，其通过Score Distillation Sampling利用文本到图像扩散模型指导三维表示优化，从而生成符合文本描述的三维对象\citep{poole2022dreamfusion}。这种方法说明，即使没有大规模三维标注数据，也可以借助二维生成模型中学习到的物体外观先验来约束三维生成。

与文本到三维不同，本文任务的输入是鞋履草图，因此更接近图像条件三维生成。Zero-1-to-3提出从单张图像预测目标物体新视角的扩散模型，将相机视角变化作为条件，生成多视角图像以支持后续三维重建\citep{liu2023zero123}。该类方法的核心思想是：先补全目标物体在不同视角下的外观，再根据多视角信息构建三维形状。这与本文关注的单视图到三维和三视图融合问题高度相关。

One-2-3-45进一步将单图像新视角生成和三维重建流程连接起来，目标是在较短时间内从单张图像得到三维网格\citep{liu2023one2345}。后续One-2-3-45++引入更一致的多视图生成和三维扩散机制，以提升生成视图之间的一致性和三维结果质量\citep{liu2023one2345pp}。这些工作表明，单图像到三维资产生成正在从逐样本优化转向更高效、更自动化的流程。

扩散先验方法对本文的启发在于：鞋履草图本身信息不足，需要模型根据学习到的外观和形状先验补全三维结构；但其局限也很明显。扩散模型生成的是概率意义上合理的对象，而不一定严格遵守输入草图的所有设计细节，更不能保证隐藏内部结构的工程真实性。因此，本文将扩散先验视为生成外观级三维模型的有效基础，而不是制造级结构重建的充分条件。

\section{多视图生成与前馈式三维重建}

多视图生成方法试图通过生成或利用多个视角的图像来缓解单视图信息不足的问题。若将第$i$个视角的图像记为$I_i$，相机投影或渲染函数记为$\Pi_i(\cdot)$，多视图重建可以抽象为同时满足多个视角观测的一致性优化：
\begin{equation}
  \hat{S} = \arg\min_S
  \sum_{i=1}^{N} \mathcal{D}\left(\Pi_i(S), I_i\right)
  + \lambda\mathcal{R}(S),
  \label{eq:multiview-consistency}
\end{equation}
其中$\mathcal{D}$衡量渲染结果与输入视图之间的差异，$\mathcal{R}$表示几何正则项。Wonder3D通过跨域扩散生成一致的多视角图像和法线信息，以提升单图像到三维重建的几何一致性\citep{long2023wonder3d}。这类方法对于鞋履设计任务具有直接意义，因为鞋履设计中常见的侧视图、俯视图和底视图本身就是多视图表达形式。若模型能够稳定利用这些视图，理论上可以更好地约束鞋头、鞋底、鞋口和后跟等结构。

近年来，前馈式三维重建方法进一步提升了推理效率。LGM提出大规模多视图高斯模型，用于从输入图像快速生成三维内容\citep{tang2024lgm}。TripoSR则面向快速单图像三维物体重建，强调较短推理时间和易用性\citep{tochilkin2024triposr}。这些方法反映出图像到三维任务正在向实时或近实时应用靠近，也为本文构建本地演示系统提供了技术背景。

Stable Fast 3D是本文早期实验中采用的稳定基线。该方法面向快速三维网格重建，并支持UV展开和光照分离，适合从单张图像生成可查看的三维资产\citep{boss2024sf3d}。在本文项目实验中，SF3D能够稳定生成可加载的鞋履GLB模型，因此被保留为基线和回退后端。但人工观察也表明，SF3D在鞋履内部结构上存在明显不足，尤其在单视图输入条件下难以生成合理的鞋口和内部空间。

Hunyuan3D-2是本文当前选择的主要生成后端相关工作。该系列模型面向高分辨率纹理三维资产生成，包含形状生成和纹理生成等模块\citep{hunyuan3d2025hunyuan3d2}。在本项目本地测试中，Hunyuan3D-2mini的形状生成能够顺利运行，并在鞋履样例上获得比SF3D更好的外观效果。因此，本文将Hunyuan3D-2mini作为当前默认形状生成后端，并将Hunyuan3D系列的多视图能力作为后续三视图输入优化的重要方向。

总体而言，多视图和前馈式三维重建方法为本文提供了两类重要启发：一方面，多视图输入能够更好地符合鞋履设计表达习惯；另一方面，前馈式模型更适合毕业设计演示所需的交互效率和本地可复现性。但在有限显存和开源模型能力边界下，三视图融合仍需要分阶段推进，不能直接假设已经达到工业级建模精度。

\section{草图条件控制与领域适配}

鞋履设计草图与自然图像存在显著差异。自然图像通常包含连续颜色、真实材质、阴影和背景上下文，而设计草图主要由轮廓线、结构线和少量明暗信息组成。若直接将草图输入图像到三维模型，模型可能难以理解线条对应的真实体积关系，生成结果容易出现扁平、缺失或结构漂移。因此，在草图输入和三维生成模型之间引入领域适配模块是本文系统路线的重要组成部分。

ControlNet为这一问题提供了可行工具。该方法在预训练扩散模型的基础上增加可训练的条件控制分支，使模型能够根据边缘、草图、姿态、深度等控制图生成图像\citep{zhang2023controlnet}。从条件生成角度看，ControlNet可以理解为在式\eqref{eq:diffusion-loss}中的条件$c$之外额外引入控制图$h$，从而使噪声预测网络变为$\epsilon_\theta(x_t,t,c,h)$。对于本文任务，ControlNet可以使用Canny边缘图或Scribble风格控制图，将鞋履设计线稿转换为更接近产品渲染图的中间图像，从而降低下游图像到三维模型的输入域差异。图\ref{fig:sketch-domain-adaptation}给出了本文采用的草图域适配思路。

\begin{figure}[htbp]
  \centering
  \includegraphics[width=0.95\textwidth]{related_work/sketch_domain_adaptation.png}
  \caption{草图域适配与图像到三维生成流程示意}
  \label{fig:sketch-domain-adaptation}
\end{figure}

除通用ControlNet方法外，草图到三维重建也是一个独立研究方向。近期关于多草图三维重建的工作尝试从多个草图中恢复三维场景结构，说明线稿输入可以作为三维恢复的有效约束\citep{sketch3d2025reconstruction}。不过，草图三维重建通常需要处理线条歧义、视角对应和几何闭合等问题，且面向通用场景或特定对象时方法设计差异较大。对于鞋履设计而言，草图中的线条既可能表示真实边界，也可能表示缝线、装饰、材质分割或设计标注，这进一步增加了从草图到几何结构的映射难度。

本文采用的策略是将草图域适配与开源三维生成后端结合：首先保留原始草图作为输入证据，然后生成边缘或涂鸦控制图，再通过ControlNet得到更模型友好的鞋履渲染图，最后送入Hunyuan3D或SF3D生成三维模型。该策略的优点是工程实现可控，能够在本地硬件上运行，并且便于保存中间结果用于论文分析；局限是ControlNet生成的渲染图可能引入与原始设计不一致的细节，因此系统需要在实验中同时展示原始草图、控制图、渲染图和最终三维模型，避免只评价最终外观而忽略设计保持程度。

\clearpage
\section{本章小结}

本章回顾了与鞋履设计草图到三维模型重建相关的主要研究方向。单图像三维重建方法为从二维视觉特征预测三维形状提供了基础，但对草图输入和隐藏结构推断仍存在不足。基于扩散先验的三维生成方法能够利用大规模二维生成模型中的外观先验，提升从有限输入生成三维内容的能力，但其结果更多体现视觉合理性而非严格工程正确性。多视图生成和前馈式三维重建方法提升了三维生成效率，也为三视图鞋履草图输入提供了技术方向。ControlNet和草图条件控制方法则为缓解设计草图与自然图像之间的域差异提供了可实施方案。

基于上述分析，本文选择“Hunyuan3D-2mini主后端、SF3D基线、ControlNet草图域适配”的系统路线。该路线既利用了当前开源图像到三维模型的生成能力，又正视了本地显存、草图输入和鞋履内部结构推断的限制。后续章节将在此基础上介绍本文系统的具体方法设计、工程实现和实验分析。
